\scanpage{155}
\deleted{Lorsque \unsure{cependant} dans $C$ les produits fibrés existent,
alors il n'y a pas de différence entre épim. eff. et épim. strict, il n'y a plus que
quatre (au lieu de cinq) variantes d'épimorphicité pour des familles, et ces
\deleted{variantes} propriétés se reconnaissent déjà sur \deleted{\unsure{R}}
le crible $R$ (et y prennent le $=$ nom).

\uline{Prop.} 0 Les \deleted{cinq} notions \deleted{\unsure{envisagées}}
\inserted{``épim.'', ``épim \unsure{stricts|stricte}'' \deleted{\unsure{épim}}}
sont \deleted{\unsure{les propri}} stables par agrandissement de la famille;
de $=$ pour les trois autres variantes, \uuline{\unsure{si on n'ajoute}}
que des \unsure{hypothèses} favorables.

Soit $G$ une partie de $\operatorname{Ob}C$, $\underline G$ la sous-catégorie pleine
engendrée. \deleted{Pour tout}

\uline{Prop.} 1 Conditions équivalentes

$a)$ Pour tout $X\in\operatorname{Ob}C$, la famille des flèches de but $X$,
sources $\in G$, est épimorphique (resp. épimorphique stricte).

$b)$ Pour tout $X\in\operatorname{Ob}C$, considérons
\diagram{\begin{tikzcd}\underline G_{/X}\arrow[r,"i"]&C\end{tikzcd}}
et le système inductif évident
\[
\varinjlim_{T\in\underline G_{/X}}T\longrightarrow X,
\]
[dans \inserted{i.e. en sens inverse} les foncteurs \unsure{covariants en $Y$}
\inserted{cov.}], ce système est épimorphique \deleted{\unsure{strict}}
(resp. isomorphique) [i.e. l'hom de foncteurs co. cov. en $Y$ est un mono,
resp. un isom].

\uline{Dém} Le cas non resp. est une tautologie, regardons le cas resp.
La condition $a)$ signifie que, si $u_i:G_i\longrightarrow X$ sont les différents
objets de $\underline G_{/X}$, alors $\forall Y$
$(i)$ $h_Y(X)\longrightarrow\prod_i h_Y(G_i)$ injectif et $(ii)$ Pour des
$\unsure{\xi_i|\zeta_i}\in h_Y(G_i)$, tels}
\margin{154}
